On the exactly solvable 2D HP lattice with 16 residues, a conditional autoregressive designer trained on uniquely folding sequences generalises to held-out target conformations: it designs \rhval{main/conditional-ar-designer/hp16/succ_k100/mean:3} of them at an oracle budget of $K=100$, far above the registered annealing baseline (\rhval{main/simulated-annealing/hp16/succ_k100/mean:3}) and above the contact heuristic (\rhval{main/contact-heuristic/hp16/succ_k100/mean:3}). Conditioning on the target is essential (unconditional model: \rhval{abl_components/unconditional-ar/hp16/succ_k100/mean:3}). The result is mixed in three respects: the hand-made contact heuristic is at least as good at $K=1$ with no learning at all; reversal augmentation is a no-op at full data (the reversed pairs are already in the training set, so the registered ablation measures training length) and, in a post-hoc follow-up, helps $K=1$ success only at small data fractions, while more scored sequences raise $K=1$ success but barely change $K=100$; and annealing with a graded degeneracy penalty, started from the heuristic sequence, is better than the AR model at $K=1000$ (\rhval{extra_sa_log/sa-graded-objective-heuristic-init/hp16/succ_k1000/mean:3} against \rhval{main/conditional-ar-designer/hp16/succ_k1000/mean:3}). The honest summary is that the learned designer's advantage is a small-to-medium-budget advantage over search on this toy, and that the comparison is more sensitive to the baseline's objective and initialisation than to the choice of a neural versus a search-based designer.
